\documentclass[conference]{IEEEtran}
\IEEEoverridecommandlockouts

\usepackage{cite}
\usepackage{amsmath,amssymb,amsfonts}
\usepackage{graphicx}
\usepackage{booktabs}
\usepackage{array}
\usepackage{textcomp}
\usepackage{xcolor}
\usepackage{url}
\def\UrlFont{\rmfamily}

\begin{document}

\title{RAGnosis: A Deterministic-Core, Retrieval-Augmented Platform for
Medical Report Analysis with Clinician-in-the-Loop Validation}

\author{\IEEEauthorblockN{[First Author], [Second Author], [Third Author]}
\IEEEauthorblockA{\textit{Department of [Department Name]} \\
\textit{[Institution Name]}\\
{[City], India} \\
\{[email1], [email2], [email3]\}@[institution-domain]}
\and
\IEEEauthorblockN{Guide: [Guide Name]}
\IEEEauthorblockA{\textit{Department of [Department Name]} \\
\textit{[Institution Name]}\\
{[City], India} \\
{[guide-email]@[institution-domain]}}}

\maketitle

\begin{abstract}
Patients in India routinely receive laboratory reports as dense printed
tables or photographs of printed tables, and most lack the background to
interpret values that fall outside reference ranges. General-purpose
large language models (LLMs) can explain such reports, but when asked to
read them directly they may transcribe values incorrectly---an
unacceptable failure mode for medical numbers. We present RAGnosis, a
retrieval-augmented platform for medical report analysis integrated with
hospital management (patient, doctor, and receptionist portals). Its
central design decision is a strict separation of concerns: a
\emph{deterministic analysis core}---PDF text extraction, dual-pass
optical character recognition (OCR) for photographed reports, rule-based
metric extraction against canonical reference ranges, a transparent
health score, and rule-scored multi-label panel classification---performs
every numeric operation, while generative models are used only to
explain values that have already been extracted and verified. A
report-grounded dialogue module answers patient questions primarily from
the structured extracted metrics, intercepting diagnosis requests and
pivoting to the patient's own flagged values. A clinician-in-the-loop
workflow lets doctors correct extracted values after the fact; corrections
recompute all derived quantities deterministically and are audit-logged.
For the mobile client, the sentence-embedding model used for retrieval
is exported to ONNX and int8-quantized for on-device inference. We
describe the system architecture, the deterministic pipeline, the
guardrailed dialogue design, and a rigorous evaluation protocol for
validation on a corpus of Indian laboratory reports that is currently
being assembled; large-scale accuracy results are therefore reported as
pending rather than preliminary.
\end{abstract}

\begin{IEEEkeywords}
retrieval-augmented generation, medical report analysis, OCR,
deterministic extraction, multi-label classification, clinician-in-the-loop,
quantized inference, health informatics
\end{IEEEkeywords}

\section{Introduction}
Laboratory reports are the most common artefact of outpatient medicine,
yet their presentation---grids of analytes, values, units, and reference
intervals---assumes clinical training. A patient who receives a complete
blood count (CBC) or lipid profile typically learns whether anything is
``wrong'' only at a follow-up consultation days later. The intervening
anxiety, and the temptation to paste the report into a general-purpose
chatbot, are both understandable. The risk is that a generative model
asked to read a photographed report end-to-end may silently misread a
digit, transpose two rows, or invent a value that was never measured;
in a medical context, a single wrong number is not a rounding error but
misinformation about a person's body.

This paper describes RAGnosis, a platform built around the opposite
commitment: \emph{no number shown to a patient is ever produced by a
generative model}. Every value is extracted by deterministic,
auditable rules from the report itself; every status (Normal, Low,
High) is decided by comparison against a canonical reference range;
every derived quantity is a closed-form function of extracted values.
Generative modelling is confined to two roles where its failure modes
are benign: drafting plain-language explanations of already-verified
values, and conversational answers that are grounded in the structured
extraction and cite it.

RAGnosis couples this analysis pipeline with the operational context in
which reports actually circulate---a hospital workflow with patient,
doctor, and receptionist portals, appointment scheduling, prescriptions,
and medicine reminders---and with a clinician-in-the-loop validation
workflow in which doctors review and correct machine extractions of
their own patients' reports. The contributions of this work are:

\begin{itemize}
\item A deterministic extraction pipeline for Indian laboratory
reports, covering digital PDFs and photographed reports via a dual-pass
OCR strategy that selects between plain and preprocessed readings by a
text-quality score (Section~IV).
\item A rule-scored \emph{multi-label} classification scheme over 13
report panels, combining keyword evidence with extracted-parameter
evidence, which reflects the reality that one physical report often
contains several panels (Section~IV).
\item A report-grounded dialogue design in which structured extracted
metrics---not raw OCR text---are the primary context, with guardrails
that intercept diagnosis requests and pivot to the patient's flagged
values (Section~V).
\item A clinician-in-the-loop validation workflow with deterministic
recomputation and audit logging of every correction (Section~VI).
\item An on-device inference path for the mobile client based on an
int8-quantized ONNX export of the retrieval embedding model
(Section~VII), and an evaluation protocol for a forthcoming corpus of
Indian laboratory reports (Section~VIII).
\end{itemize}

\section{Related Work}
Retrieval-augmented generation (RAG), introduced by Lewis \emph{et
al.}~\cite{lewis2020rag}, conditions a sequence-to-sequence model on
passages retrieved from an external corpus, improving factuality on
knowledge-intensive tasks. We adopt the retrieval pattern but invert
the usual trust relationship: retrieved passages serve only as
background education, while the authoritative context for a patient's
question is a structured record produced by deterministic extraction.

Deep learning has been applied broadly in healthcare~\cite{esteva2019guide},
including report interpretation; however, end-to-end neural reading of
documents couples transcription and interpretation in a way that makes
individual values hard to audit. Our pipeline instead composes
established, inspectable components: the Tesseract OCR
engine~\cite{smith2007tesseract} for photographed reports and
pdfplumber~\cite{pdfplumber} for digital PDFs. For retrieval we use
sentence embeddings in the style of Sentence-BERT~\cite{reimers2019sbert}
built on BERT~\cite{devlin2019bert}, specifically the distilled MiniLM
model of Wang \emph{et al.}~\cite{wang2020minilm}, indexed with
FAISS~\cite{johnson2019faiss}. Model compression via quantization has a
long history in efficient inference~\cite{jacob2018quantization};
we apply post-training int8 dynamic quantization to the embedding model
for on-device use (Section~VII). The generative component, when
configured, is LLaMA~3.3-70B~\cite{llama3herd2024} served through the
Groq LPU inference platform~\cite{abts2020groq}; the system is designed
so that no functionality depends on its availability.

Clinician-in-the-loop and human-in-the-loop machine learning more
generally recognise that expert review is not merely quality assurance
but a data-generating process: corrections become labelled examples.
Our validation workflow (Section~VI) follows this pattern while keeping
the correction semantics deterministic---a corrected value is
reclassified by the same rules as an original extraction, so the audit
trail explains itself.

\section{System Architecture}
RAGnosis is a three-tier system (Fig.~\ref{fig:arch}). The presentation
tier is a React~18 single-page application (built with Vite) that serves
three role-specific portals---patient, doctor, and
receptionist---alongside a React Native (Expo) mobile client. The
application tier is a Flask~3 API served by Gunicorn, organised into
modules for authentication, reports, chat, hospital operations, and
reminders, and backed by MongoDB Atlas. The intelligence tier is the
analysis pipeline of Section~IV plus the dialogue module of
Section~V. Optional components (the LLM provider, the embedding model,
and the OCR engine binary) are imported lazily: the platform boots and
serves with none of them installed, degrading to deterministic
behaviour rather than failing.

\begin{figure}[t]
\centering
\small
\begin{tabular}{|p{0.92\columnwidth}|}
\hline
\textbf{Presentation tier} \\
React 18 + Vite web portals (patient / doctor / receptionist);
Expo (React Native) mobile client \\
\hline
\textbf{Application tier (Flask 3 + Gunicorn)} \\
Auth (PyJWT, bcrypt) $\cdot$ Reports $\cdot$ Chat $\cdot$ Hospital
(appointments, prescriptions) $\cdot$ Reminders \\
\hline
\textbf{Deterministic analysis core} \\
pdfplumber / dual-pass OCR $\rightarrow$ metric extraction $\rightarrow$
multi-label classification $\rightarrow$ alerts $\rightarrow$ health
score $\rightarrow$ summary \\
\hline
\textbf{Dialogue tier} \\
Structured report context $\cdot$ optional LLaMA-3.3-70B (Groq)
$\cdot$ deterministic report-data answerer $\cdot$ FAISS + MiniLM
retrieval (background) \\
\hline
\textbf{Data tier} \\
MongoDB Atlas: users, reports, appointments, prescriptions,
reminders, chat sessions, validation corrections, audit log \\
\hline
\end{tabular}
\caption{RAGnosis tier architecture. Data flows upward from extraction
to presentation; no tier below the dialogue tier contains a generative
model, and no numeric result originates in the dialogue tier.}
\label{fig:arch}
\end{figure}

Security is enforced at the application tier. Authentication uses
short-lived JWT access tokens (30 minutes) with rotating refresh
tokens; passwords are hashed with bcrypt. Every route is guarded by
role-based access control and per-object ownership checks---a patient
can only ever read objects whose \texttt{patient\_id} equals their own
identity, and attempts to cross that boundary are written to an audit
log. Uploads are validated by magic bytes and size limits, login and
password-reset endpoints are rate-limited, and password-reset tokens
are single-use, short-lived, and stored only as SHA-256 hashes.
Deduplication of reports is by content hash per patient, so re-uploads
return the existing analysis instead of duplicating storage.

\section{Deterministic Extraction and Multi-Label Classification}

\subsection{Text Acquisition}
Digital PDFs are read with pdfplumber~\cite{pdfplumber}, extracting both
flowing text and table cell contents. Reports that exist only as
photographs or scans---the common case for patients---are processed by
an OCR path built on OpenCV pre-processing and
Tesseract~\cite{smith2007tesseract}. Photographed tables defeat naive
OCR in two distinct ways: glare, skew, and low contrast corrupt glyphs,
while aggressive binarisation can shred the thin strokes of an
otherwise clean image. Our pipeline therefore runs a \emph{dual-pass}
read: a plain grayscale pass first, and, only if its text quality is
poor, a second pass over an adaptively thresholded, deskewed image.
Each candidate text is scored by a deterministic quality function
(the count of word-like tokens), and the higher-scoring text is kept.
The plain-first order matters: in practice the pre-processing chain
that rescues a difficult photograph can destroy a crisp screenshot,
and the quality pick makes the choice per document rather than by
fiat. Both passes use Tesseract's uniform-block page segmentation,
which preserves row order in tabular material far better than fully
automatic segmentation.

\subsection{Metric Extraction}
Extraction is a rule system over a canonical analyte table. Each entry
defines a parameter (e.g.\ haemoglobin, platelets, LDL cholesterol),
its synonyms as printed by different laboratories, a canonical unit, a
unit-conversion factor where laboratories differ, and a canonical
reference range. Three complementary passes run over the text:
(i)~a \emph{line pass} matching ``name value unit'' patterns on a single
line; (ii)~a \emph{sequence pass} for OCR output in which a table has
been dumped column-wise, pairing the ordered list of parameter names
with the ordered list of values that follows it; and (iii)~a
\emph{dense-table fallback} (per the project specification) that
line-scans crowded layouts which defeat the first two passes. Values
are normalised to canonical units before comparison. Each extracted
metric is then assigned a status by direct comparison with the
canonical range:
\begin{equation}
\text{status}(v) =
\begin{cases}
\text{Low}, & v < \text{lo} \\
\text{Normal}, & \text{lo} \le v \le \text{hi} \\
\text{High}, & v > \text{hi}
\end{cases}
\end{equation}
Statuses deliberately follow the canonical ranges rather than ranges
printed on the report, because printed ranges vary by laboratory and by
assay; the canonical table is versioned with the code and is the same
for every patient, which is what makes clinician corrections
(Section~VI) comparable across reports.

\subsection{Alerts and Health Score}
A small set of critical-value rules (for example, severely low platelet
counts) raises alerts that are surfaced on the report and to the
patient's doctor. The health score is intentionally transparent:
\begin{equation}
\text{score} = \frac{\#\{\text{metrics with status Normal}\}}
{\#\{\text{metrics extracted}\}} \times 100
\end{equation}
We make no claim that this ratio is a clinical risk measure; it is a
communication device whose formula is disclosed to the user in the
interface, and its simplicity is the point---a patient can verify it by
counting.

\subsection{Multi-Label Panel Classification}
A physical report frequently bundles several panels: one sheet may
carry a CBC, a lipid profile, and a fasting glucose. A single-label
classifier is therefore the wrong abstraction. RAGnosis scores every
one of 13 panel categories---CBC, Lipid Profile, Kidney Function,
Liver Function, Thyroid, Diabetes, Urine, Radiology, Cardiac Markers,
Electrolytes, Vitamins, Hormones, and General---by two evidence
channels (Table~\ref{tab:classify}):

\begin{table}[t]
\centering
\small
\caption{Multi-label classification evidence. A category is assigned
when its combined evidence score reaches a fixed threshold; the
highest-scoring category is retained as the primary label.}
\label{tab:classify}
\begin{tabular}{@{}ll@{}}
\toprule
\textbf{Channel} & \textbf{Evidence} \\
\midrule
Keyword & Panel names and characteristic phrases in the \\
        & report text (e.g.\ ``LIPID PROFILE'', ``HbA1c'') \\
Parameter & Extracted metrics that belong to the panel \\
        & (e.g.\ haemoglobin, platelets $\rightarrow$ CBC) \\
\bottomrule
\end{tabular}
\end{table}

Each matched keyword and each panel-consistent extracted parameter adds
a weighted contribution; the normalised score is reported as an
evidence-strength confidence per label. All labels at or above the
threshold are stored on the report, so downstream consumers---doctor
suggestions, summary focus, and the doctor's correction workflow---can act
on the full panel set, while the highest-scoring label remains
available as the primary category for list views and filtering. Because
parameter evidence derives from extraction, classification quality is
bounded by extraction quality, and both are measured separately in the
evaluation protocol (Section~VIII).

\subsection{Summarisation}
Summaries are drafted from the extracted metrics, never from raw text:
an extractive template composes the in-range and flagged values, the
score, and the panel labels; when the LLM is configured it may rewrite
the same facts in plainer language under a system prompt that forbids
introducing new numbers, and the extractive version remains the
fallback. The summary carries an educational-only disclaimer wherever
it is displayed.

\section{Report-Grounded Dialogue with Guardrails}
The dialogue module answers a patient's questions about their own
report. Its defining choice is the context hierarchy. The \emph{primary}
context is the structured extraction: parameter names, values, units,
canonical ranges, statuses, the health score, and the summary. The
\emph{secondary} context is retrieval over a small curated knowledge
base of educational passages, embedded with
MiniLM~\cite{wang2020minilm} in the Sentence-BERT
style~\cite{reimers2019sbert} and indexed with
FAISS~\cite{johnson2019faiss}; retrieval supplies background explanation
only and is never treated as a source of patient facts. The raw OCR
text is used, if at all, only as a truncated tail excerpt, because for
photographed reports it is a column-wise dump in which values are
detached from their labels---precisely the failure mode the structured
pipeline exists to repair.

Answers are produced by the first available of three answerers:
(1)~the LLM (LLaMA~3.3-70B~\cite{llama3herd2024} via
Groq~\cite{abts2020groq}), given the structured context and any
retrieved passages under rules that require quoting values exactly and
forbid diagnosis; (2)~a \emph{deterministic report-data answerer} that
answers from the metrics directly---per-parameter value, range, and
status; the flagged-value list; the health score; a full listing; and
(3)~curated educational passages, personalised with the patient's own
value for that parameter when one exists. The deterministic answerer is
not a stub: it is the default in deployments without an LLM key, and
it guarantees that the most safety-critical questions (``which of my
values are out of range?'') are answered by code whose behaviour can be
tested exhaustively.

Guardrails operate before any answerer. Diagnosis-seeking questions
(``do I have \ldots?'') are intercepted and answered with an explanation
of what the system can do, pivoted to the patient's own flagged values
so the response remains useful rather than a refusal wall.
Prescription and dosage requests are declined. When no report is
attached to a session, the assistant does not fabricate one; it states
that no report data is available and answers only general educational
questions. Every answer carries the educational, non-diagnostic
disclaimer described in Section~IX.

\section{Clinician-in-the-Loop Validation}
Extraction from photographed reports will sometimes be wrong, and a
system that shows patients numbers needs a correction path that is
stronger than a feedback button. Reports in RAGnosis are visible to
the patient immediately---visibility is never gated on clinician
action---and correction is layered on top as an optional, auditable
oversight path. From the patient record, the treating doctor sees the
extracted values of that patient's reports and can correct any value
that the pipeline read wrongly. There is deliberately no confirm or
approve action: silence means the values stand as extracted, and only
actual clinician edits enter the record.

Corrections are applied through the deterministic core, not around it:
a corrected value is re-compared against the canonical range to
recompute its status, and the health score, alerts and plain-language
summary are recomputed. Every correction writes an
audit record (who, when, original value, corrected value), and the
report thereafter displays a ``Values corrected by Dr.\ [name]'' note
to the patient, distinguishing clinician-corrected numbers from
uncorrected extractions. The audit trail doubles as the annotation substrate for
the evaluation corpus of Section~VIII: a corrected field is a labelled
example with a known provenance, produced as a side effect of clinical
use rather than a separate labelling campaign. Access control is
unchanged by the workflow---a doctor can correct only reports of
patients with whom an appointment relationship exists, enforced by the same
ownership checks as every other route.

\section{On-Device Quantized Inference}
The mobile client must answer retrieval-backed questions responsively
and, ideally, with some capability when connectivity is poor. Server
round-trips for embedding are the avoidable cost: embedding is needed
for every question, is small compared with generation, and is
privacy-sensitive because the query is derived from a medical report.
We therefore export the retrieval encoder (all-MiniLM-L6-v2, a 6-layer
distilled transformer~\cite{wang2020minilm}) to ONNX and apply int8
\emph{dynamic} quantization---weights stored in int8, activations
quantized at runtime---following the integer-inference approach of
Jacob \emph{et al.}~\cite{jacob2018quantization}, and run it with ONNX
Runtime~\cite{onnxruntime} inside the Expo/React Native client. The
deterministic answerer also runs in this path conceptually unchanged:
it consumes structured metrics, which are computed server-side at
ingestion and cached on the device with the report.

Because the mobile export and its device benchmarks are being finalised
in parallel with this draft, we report the measurement plan rather than
numbers we cannot yet stand behind (Table~\ref{tab:ondevice}). Parity
is defined as agreement between on-device and server embeddings:
cosine similarity on a fixed probe set, and retrieval overlap (whether
the top-$k$ passages returned on-device match the server's) on
representative patient questions.

\begin{table}[t]
\centering
\small
\caption{On-device embedding model (all-MiniLM-L6-v2), measured on CPU
with the export pipeline of Sec.~VII. Per-channel dynamic int8
quantization is essential: per-tensor scales measured only 0.966 mean
cosine parity, against 0.993 per-channel at the same size.}
\label{tab:ondevice}
\begin{tabular}{@{}lll@{}}
\toprule
\textbf{Quantity} & \textbf{FP32} & \textbf{INT8} \\
\midrule
Model size on disk        & 90.4 MB & 23.0 MB \\
Embedding latency (batch 1, CPU) & 5.6 ms & 4.0 ms \\
Cosine parity vs.\ server  & 1.000 (ref.) & 0.993 \\
Retrieval top-1 agreement  & 1.00 & 1.00 \\
\bottomrule
\end{tabular}
\end{table}

\section{Evaluation Methodology and Preliminary Verification}

\subsection{Corpus and Annotation Protocol}
The evaluation corpus is being assembled from partner laboratories in
India and is deliberately \emph{not} part of this draft; we describe
the protocol so that the eventual numbers are interpretable.
Collection targets several hundred reports across the 13 panels of
Section~IV, with a deliberate majority of photographed (rather than
digitally native) reports, since photographs dominate real patient
uploads. Each report is de-identified before annotation: names,
identifiers, contact details, and laboratory accession numbers are
removed, and images are cropped to the report body where feasible.
Two clinicians independently annotate, per report: the set of panel
labels; and for each parameter, the printed value and unit transcribed
verbatim. Disagreements are adjudicated by a third clinician. The
result is a gold standard against which extraction is scored field by
field rather than by document-level impression.

\subsection{Metrics}
We will report: (i)~\emph{extraction quality}---precision, recall, and
F1 per parameter and micro/macro-averaged, where a field counts as
correct only if value (within rounding tolerance) and unit both match;
(ii)~\emph{status agreement}---the fraction of extracted fields whose
Normal/Low/High status matches the status implied by the gold value
under the canonical ranges; (iii)~\emph{label quality}---per-label and
micro/macro F1 for the multi-label classifier, plus primary-label
accuracy; (iv)~\emph{clinician agreement}---inter-annotator agreement
(Cohen's $\kappa$) on labels and on transcription, to establish the
human ceiling; and (v)~\emph{dialogue grounding}---on a fixed question
set per report, the fraction of answers whose quoted values exactly
match the structured extraction (for the deterministic answerer this is
correct by construction and serves as a regression check; for the LLM
path it is the substantive measure).

\subsection{Preliminary System Verification}
Pending the corpus, the implemented system has been verified
functionally, end to end, on a local deployment with seeded accounts
for all three roles: an automated suite of 65 API-level checks covering
authentication and session handling, report upload and deduplication,
OCR and PDF extraction, chat grounding, booking conflicts, prescriptions,
and reminders passed in full, as did 37 browser-driven checks across the
patient, doctor, and receptionist portals including a mobile-viewport
pass. The deterministic components are pure functions and are exercised
with fixed fixtures (including a photographed CBC whose column-wise OCR
dump defeats a line-only extractor but is recovered by the sequence
pass). These are correctness checks of behaviour, \emph{not} accuracy
measurements, and we do not present them as such; the accuracy claims
of this system stand or fall on the corpus evaluation above.

\section{Limitations and Ethics}
Several limitations are intrinsic to the design. Extraction quality
bounds everything downstream: a value missed by all three extraction
passes is invisible to classification, scoring, and dialogue alike, and
handwritten annotations on reports are out of scope. Canonical reference
ranges are a simplification---true ranges vary with assay, age, sex,
and pregnancy---and the system mitigates rather than resolves this by
using one disclosed table and routing doubtful extractions to
clinicians. The health score is a communication device, not a risk
score, and its presentation says so. The LLM path, when enabled, can
still phrase an explanation misleadingly even when its numbers are
grounded; the deterministic answerer exists partly as a check on this.

Ethically, the system is educational and non-diagnostic by design: it
never states or implies a diagnosis, declines prescription requests,
and attaches disclaimers to summaries and answers. Patient data is
protected by role-based access, ownership checks, hashed credentials
and reset tokens, rate limits, and audit logging; the corpus protocol
requires de-identification before annotation and institutional ethics
clearance before any report is used for evaluation. Correction
notes are shown honestly---only reports a clinician actually
corrected carry them---because a false signal of verification would be worse than
none.

\section{Conclusion}
RAGnosis demonstrates that a patient-facing medical report platform can
be built so that generative models never originate a number: a
deterministic core extracts, classifies, and scores; a grounded
dialogue layer explains; clinicians validate through a workflow whose
corrections recompute deterministically and accumulate as labelled
data; and a quantized on-device encoder keeps retrieval responsive on
mobile. The immediate next step is the corpus evaluation of
Section~VIII, whose protocol is fixed in advance so that the results,
whatever they are, will be interpretable. We believe the
deterministic-core pattern generalises beyond laboratory reports to
any domain where an LLM's fluency is welcome but its arithmetic is
not admissible evidence.

\bibliographystyle{IEEEtran}
\bibliography{references}

\end{document}
